% --- PAGE 17: HALF-TITLE LITERATURE SURVEY ---
\vspace*{\fill}
\begin{center}
  {\fontsize{18pt}{24pt}\selectfont\bfseries LITERATURE SURVEY\par}
\end{center}
\vspace*{\fill}
\newpage

% --- PAGE 18: CHAPTER 2 LITERATURE SURVEY ---
\begin{center}
  {\fontsize{16pt}{20pt}\selectfont\bfseries CHAPTER 2\par}
  \vspace{0.3cm}
  {\fontsize{16pt}{20pt}\selectfont\bfseries LITERATURE SURVEY\par}
\end{center}
\vspace{0.25cm}

{\fontsize{11pt}{15.5pt}\selectfont
\textbf{``Development and Classroom Evaluation of a Static-Assessment Programming Practice Platform''} by Jeon, J., Jang, J., and Jung, S. K. (2025) [1] -- The authors investigated web-based programming platforms (Code-Assist-Lab) that utilize static code evaluation paradigms---specifically abstract syntax tree (AST) inspection, cyclomatic complexity calculations via Radon, and style linting via Flake8. By circumventing server-side dynamic code execution, the platform mitigates malicious process execution hazards and reduces cloud infrastructure costs. However, the architecture is strictly confined to static grading and cannot verify runtime edge-case compliance or memory/CPU resource bounds. Furthermore, it lacks collaborative peer hierarchies and longitudinal student tracking across institutional departments.

\vspace{0.2cm}
\textbf{``Integration of Multiple Sources to Anticipate Student Performance Using Learning Analytics''} by Moreno-Marcos, P. M. et al. (2025) [2] -- This research explored predictive learning analytics in higher education by consolidating heterogeneous data streams, including Learning Management System (LMS) clickstream telemetry, SPOC video viewing frequencies, and historical grade records. Through multi-source feature extraction, predictive models successfully forecasted student examination trajectories. Nevertheless, a major research gap identified was the absence of real-time intervention mechanisms; the system operated in an open-loop post-hoc reporting model lacking closed-loop remedial assignment dispatch, automated peer alerts, or synchronized hands-on coding environments.

\vspace{0.2cm}
\textbf{``Shaking The Silos: Impact of Sequential Live Coding on Students' Performance''} by Cho, K., Ali, S. F., Bang, S. J., and Anwar, S. (2024) [3] -- The authors presented an active classroom pedagogical intervention termed Sequential Live Coding (SLC), where students consecutively assume code-authoring responsibilities and articulate syntactic and algorithmic rationales to peers. Although empirical results indicated marked improvements in classroom engagement, the technique was confined to synchronous, isolated instructor-led sessions. It lacked asynchronous telemetry tracking across external competitive coding ecosystems (e.g., LeetCode streaks and GitHub commits) and provided no automated verification pipeline or institutional audit governance.

\par}
\newpage

% --- PAGE 19: CHAPTER 2 CONTINUATION & COMPARATIVE MATRIX ---
{\fontsize{11.5pt}{16.5pt}\selectfont
\noindent In order to contextualize the core technical challenges and delineate the boundaries of existing academic methodologies, Table~2.1 provides an empirical comparison of recent state-of-the-art literature along with their identified drawbacks and limitations.

\vspace{0.15cm}
% ==============================================================================
% TABLE 2.1: LITERATURE SURVEY COMPARATIVE SUMMARY MATRIX
% (Easily editable: modify, add, or remove rows using the structure below)
% ==============================================================================
\begin{table}[H]
  \centering
  \renewcommand{\arraystretch}{1.35}
  \fontsize{9pt}{12.5pt}\selectfont
  \begin{tabular}{|>{\centering\arraybackslash}m{0.9cm}|>{\raggedright\arraybackslash}m{3.6cm}|>{\centering\arraybackslash}m{2.3cm}|>{\centering\arraybackslash}m{1.1cm}|>{\raggedright\arraybackslash}m{5.7cm}|}
    \hline
    \rowcolor{black!10}
    \textbf{S.NO} & \centering\textbf{PAPER TITLE} & \textbf{AUTHOR(S)} & \textbf{YEAR} & \centering\textbf{DRAWBACKS \& RESEARCH LIMITATIONS} \tabularnewline
    \hline
    % --- ROW 1 ---
    1 & Development and Classroom Evaluation of a Static-Assessment Programming Practice Platform 
      & Jeon J. et al. 
      & 2026 
      & Confined to static automated grading; does not support collaborative peer learning or longitudinal tracking. \tabularnewline
    \hline
    % --- ROW 2 ---
    2 & Integration of Multiple Sources to Anticipate Student Performance Using Learning Analytics 
      & Moreno-Marcos P.M. et al. 
      & 2025 
      & Demonstrates predictive modeling but lacks active collaborative intervention and closed-loop remedial practice cycles. \tabularnewline
    \hline
    % --- ROW 3 ---
    3 & Shaking The Silos: Impact of Sequential Live Coding on Students' Performance 
      & Cho K.; Ali S.F.; Bang S.J.; Anwar S 
      & 2024 
      & Investigates live coding as an isolated teaching technique; lacks multi-source evidence tracking and automated remediation. \tabularnewline
    \hline
  \end{tabular}\par
  \vspace{10pt}
  {\centering\fontsize{10pt}{14pt}\selectfont\bfseries Table 2.1: Comparative Summary of Literature Survey and Identified Limitations\par}
\end{table}

\vspace{0.25cm}
\textbf{Summary of Literature Gaps and Proposed System Motivation:}\par
An analysis of the literature reveals a structural disconnect in academic computing infrastructure: platforms either offer isolated static grading (Jeon et al.~[1]), passive predictive analytics without active intervention (Moreno-Marcos et al.~[2]), or isolated synchronous live coding without automated continuous telemetry (Cho et al.~[3]). 

The proposed \textbf{Academic Task Governance and Code Assessment Platform} directly resolves these limitations by synthesizing: (1) an isolated multi-language compiler sandbox (C, C++, Java, Python) with Judge0 cloud failover for live runtime verification, addressing the sandboxing models formalized by Wasik et al.~[4]; (2) a formal multi-level Role-Based Access Control (RBAC) governance pipeline adhering to the authorization models established by Sandhu et al.~[5]; (3) automated LeetCode GraphQL and GitHub REST telemetry daemons providing longitudinal skill analytics; and (4) an automated 3-tier review hierarchy (Student Coordinator $\rightarrow$ Faculty Advisor $\rightarrow$ HOD) with closed-loop multi-channel alerts (Telegram, Web Push, and email) to enforce immediate remedial intervention.\par}
\newpage
